% !TEX root = ../main.tex
\begin{ouverture}[mouvement]{13}{Forces et variation de la vitesse}{Partie 2 · Mouvement et interactions}
{Vecteur vitesse\enspace·\enspace Variation de vitesse\enspace·\enspace Somme des forces\enspace·\enspace Rôle de la masse\enspace·\enspace Chronophotographie}
\objectif{Construire des vecteurs vitesse et variation de vitesse à partir d'une chronophotographie ou d'une vidéo.}
\objectif{Utiliser la relation approchée $\sum\vec{F}\approx m\,\dfrac{\Delta\vec{v}}{\Delta t}$ pour estimer une variation de vitesse, les forces étant connues.}
\objectif{Utiliser la même relation pour estimer les forces, le mouvement étant connu.}
\objectif{Expliquer le rôle de la masse ; étudier la relation à l'aide d'un programme.}
\prerequis{Référentiel, vecteur position, vecteur vitesse, principe d'inertie (classe de seconde).}
\prerequis{Forces usuelles : poids, réaction d'un support, tension d'un fil, frottements ; somme de vecteurs.}
\end{ouverture}
\essentiel
\essSection{1}{Vecteur vitesse et variation de vitesse}
\begin{definition}[Vecteur vitesse]
Dans un référentiel donné, on repère les positions $M_0$, $M_1$, $M_2$… d'un point à intervalles de temps réguliers $\Delta t$. Le vecteur
vitesse en $M_i$ est approximativement
\[\vec{v}_i\approx\frac{\overrightarrow{M_iM_{i+1}}}{\Delta t}\qquad\text{de valeur}\qquad v_i\approx\frac{M_iM_{i+1}}{\Delta t}\ \text{(en \unit{m.s^{-1}})}.\]
Il est d'autant plus proche du vecteur vitesse réel que $\Delta t$ est petit : il est alors pratiquement tangent à la trajectoire.
\end{definition}
\begin{definition}[Vecteur variation de vitesse]
Entre deux instants voisins, le vecteur variation de vitesse est $\Delta\vec{v}_i=\vec{v}_{i+1}-\vec{v}_i$. On le construit en traçant $\vec{v}_i$ et
$\vec{v}_{i+1}$ à partir d'un même point : $\Delta\vec{v}_i$ va de l'extrémité de $\vec{v}_i$ à celle de $\vec{v}_{i+1}$.
\end{definition}
\schema{13-vitesse}
\essSection{2}{Variation de vitesse et somme des forces}
\begin{loi}[Relation approchée]
Pour un système de masse $m$ modélisé par un point, entre deux instants voisins séparés de $\Delta t$ :
\[\sum\vec{F}\approx m\,\frac{\Delta\vec{v}}{\Delta t}.\]
Le vecteur $\Delta\vec{v}$ a la direction et le sens de la somme des forces. Exemple : une voiture de \qty{1000}{kg} soumise à une somme des
forces de \qty{2000}{N} dirigée vers l'avant gagne \qty{2,0}{m.s^{-1}} en \qty{1,0}{s}.
\end{loi}
\par\penalty0
\schema{13-relation}
\begin{propriete}[Rôle de la masse et principe d'inertie]
Pour une même somme des forces, la variation de vitesse est d'autant plus faible que la masse est grande : $\Delta v\approx\dfrac{\sum F\times\Delta t}{m}$.
Si les forces se compensent ($\sum\vec{F}=\vec{0}$), alors $\Delta\vec{v}=\vec{0}$ : le mouvement est rectiligne et uniforme, ou le système reste
immobile (principe d'inertie). Réciproquement, si le vecteur vitesse change, même seulement de direction, les forces ne se
compensent pas.
\end{propriete}
\begin{propriete}[Chute libre]
Un objet soumis à son seul poids (frottements négligés) vérifie $m\,\dfrac{\Delta\vec{v}}{\Delta t}=m\,\vec{g}$, soit $\Delta\vec{v}=\vec{g}\,\Delta t$ : sa vitesse varie de
\qty{0,98}{m.s^{-1}}, verticalement vers le bas, toutes les \qty{0,10}{s}, quelle que soit sa masse.
\end{propriete}
\begin{python}
dt = 0.10                             # s
y = [0.80, 0.75, 0.60, 0.36, 0.02]    # m : balle lachee, axe vers le haut
v = [(y[i + 1] - y[i]) / dt for i in range(4)]
for i in range(3):
    print('Delta v =', round(v[i + 1] - v[i], 2), 'm/s')
\end{python}
\begin{console}
Delta v = -1.0 m/s
Delta v = -0.9 m/s
Delta v = -1.0 m/s
\end{console}
\noindent\capacite{Capacité numérique} À partir des positions relevées sur une vidéo, le programme calcule les vitesses puis leurs variations :
environ $-\qty{1}{m.s^{-1}}$ toutes les \qty{0,10}{s}, soit $\dfrac{\Delta v}{\Delta t}\approx\qty{10}{m.s^{-2}}\approx g$ : la balle n'est soumise qu'à son poids.
\begin{attention}
$\Delta\vec{v}$ est une différence de vecteurs, pas de valeurs : dans un virage pris à vitesse constante, la valeur de la vitesse ne change
pas mais sa direction change, donc $\Delta\vec{v}\neq\vec{0}$ et une force s'exerce. La relation n'est qu'approchée : elle est d'autant plus
précise que $\Delta t$ est petit.
\end{attention}
\begin{savoirfaire}[Estimer une force à partir d'une chronophotographie]
\begin{enumerate}
\item Calculer $v_i$ et $v_{i+1}$ à partir des distances mesurées, converties en mètres grâce à l'échelle du document.
\item Tracer $\vec{v}_i$ et $\vec{v}_{i+1}$ depuis un même point, avec une échelle choisie (par exemple \qty{1}{cm} pour \qty{1}{m.s^{-1}}).
\item Tracer $\Delta\vec{v}_i$, qui joint leurs extrémités, et mesurer sa longueur pour obtenir $\Delta v_i$.
\item Calculer $\sum F\approx m\,\dfrac{\Delta v_i}{\Delta t}$ : la somme des forces a la direction et le sens de $\Delta\vec{v}_i$.
\end{enumerate}
Exemple : pour une balle de \qty{60}{g}, $\Delta v=\qty{0,99}{m.s^{-1}}$ vers le bas en \qty{0,10}{s}, d'où $\sum F\approx\qty{0,59}{N}$ vers le bas : c'est son poids.
\end{savoirfaire}
\begin{experience}[Exploiter une vidéo]
On filme un mouvement (chute, lancer, chariot tiré), puis un logiciel de pointage relève les positions image par image. On
construit les vecteurs vitesse et variation de vitesse, puis on compare $m\,\dfrac{\Delta\vec{v}}{\Delta t}$ à la somme des forces connues.
\end{experience}
\begin{rediger}
Préciser le système et le référentiel, faire l'inventaire des forces et un schéma, indiquer l'échelle des vecteurs, puis
écrire la relation $\sum\vec{F}\approx m\,\dfrac{\Delta\vec{v}}{\Delta t}$ avant de calculer.
\end{rediger}
\exercices

\rubrique{Découverte}
\begin{blocLivre}{ex-13-1}
\begin{exercice}[1]{Le palet sur coussin d'air}
Un palet glisse sans frottement sur une table horizontale à coussin d'air. On a enregistré la position de son centre.
\schema{13-ex01}
\begin{questions}
\item Quelle est la nature du mouvement du palet ? Justifier.
\item Calculer la valeur de la vitesse du palet au point $M_3$.
\item Représenter $\vec{v}_3$ et $\vec{v}_4$ avec l'échelle \qty{1}{cm} pour \qty{0,10}{m.s^{-1}}. Que vaut $\Delta\vec{v}_3$ ?
\item En déduire ce que l'on peut dire des forces qui s'exercent sur le palet. Les nommer.
\end{questions}
\end{exercice}
\end{blocLivre}
\begin{blocLivre}{ex-13-2}
\begin{exercice}[1]{Une balle qui tombe}
Une balle de \qty{58}{g} est lâchée sans vitesse initiale ; on relève sa position toutes les \qty{50}{ms}.
\par\smallskip\noindent\hfill\begin{tabular}{lccccccc}
\toprule
Point & $M_0$ & $M_1$ & $M_2$ & $M_3$ & $M_4$ & $M_5$ & $M_6$ \\
Distance parcourue (cm) & 0 & 1,2 & 4,9 & 11,0 & 19,6 & 30,7 & 44,1 \\
\bottomrule
\end{tabular}\hfill\null\par\smallskip
\begin{donnees}
$g=\qty{9,81}{N.kg^{-1}}$.
\end{donnees}
\begin{questions}
\item Calculer les valeurs des vitesses $v_3$ et $v_4$.
\item Représenter $\vec{v}_3$, $\vec{v}_4$ et $\Delta\vec{v}_3$ avec l'échelle \qty{1}{cm} pour \qty{1}{m.s^{-1}}.
\item Calculer $m\,\dfrac{\Delta v_3}{\Delta t}$ et le comparer au poids de la balle.
\item Conclure : quelles forces s'exercent sur la balle ?
\end{questions}
\end{exercice}
\end{blocLivre}
\begin{blocLivre}{ex-13-3}
\begin{exercice}[1]{Le rôle de la masse}
Deux chariots, de masses \qty{2,0}{kg} et \qty{10}{kg}, initialement immobiles sur un sol horizontal, sont poussés pendant \qty{0,50}{s} avec
la même somme des forces, horizontale, de valeur \qty{20}{N}.
\begin{questions}
\item Écrire la relation approchée entre la somme des forces et la variation du vecteur vitesse.
\item Calculer la variation de vitesse de chaque chariot.
\item Quel chariot est le plus difficile à mettre en mouvement ? Pourquoi ?
\item Quelle somme des forces faudrait-il appliquer au chariot de \qty{10}{kg} pour qu'il atteigne la même vitesse que l'autre ?
\end{questions}
\end{exercice}
\end{blocLivre}
\begin{blocLivre}{ex-13-4}
\begin{exercice}[1]{Un freinage}
Une voiture de \qty{1200}{kg} roulant en ligne droite à \qty{50}{km/h} s'arrête en \qty{2,5}{s}.
\begin{questions}
\item Calculer la variation de la valeur de la vitesse, en \unit{m.s^{-1}}.
\item Estimer la valeur de la somme des forces qui s'exercent sur la voiture pendant le freinage, supposée constante.
Préciser sa direction et son sens.
\item De combien la vitesse diminue-t-elle en \qty{0,10}{s} ?
\item Quelle force faudrait-il pour arrêter de la même façon un camion de \qty{12}{t} ?
\end{questions}
\end{exercice}
\end{blocLivre}
\begin{blocLivre}{ex-13-5}
\begin{exercice}[1]{Lire des vecteurs vitesse}
\schema{13-ex05}
\begin{questions}
\item Pour chaque situation, construire le vecteur $\Delta\vec{v}=\vec{v}_{i+1}-\vec{v}_i$.
\item En déduire la direction et le sens de la somme des forces.
\item Dans la situation b), quelle force s'exerce sur la balle ?
\item Dans quelle(s) situation(s) les forces se compensent-elles ?
\end{questions}
\end{exercice}
\end{blocLivre}
\rubrique{Entraînement}
\begin{blocLivre}{ex-13-6}
\begin{exercice}[2]{Une balle lancée, en Python}
Les positions d'une balle lancée ont été relevées sur une vidéo, toutes les \qty{0,10}{s}, dans un repère dont l'axe $(Oy)$ est vertical et
orienté vers le haut.
\begin{python}
dt = 0.10
x = [0, 0.30, 0.60, 0.90, 1.20, 1.50, 1.80]       # m
y = [0, 0.35, 0.60, 0.76, 0.82, 0.77, 0.63]       # m
vx = [(x[i + 1] - x[i]) / dt for i in range(6)]
vy = [(y[i + 1] - y[i]) / dt for i in range(6)]
for i in range(5):
    print('vx =', round(vx[i], 1), '| Delta vy =', round(vy[i + 1] - vy[i], 1))
\end{python}
\begin{console}
vx = 3.0 | Delta vy = -1.0
vx = 3.0 | Delta vy = -0.9
vx = 3.0 | Delta vy = -1.0
vx = 3.0 | Delta vy = -1.1
vx = 3.0 | Delta vy = -0.9
\end{console}
\begin{questions}
\item \capacite{Capacité numérique}Que calculent les listes \code{vx} et \code{vy} ?
\item Que peut-on dire de la composante horizontale de la vitesse ? Et de la variation de sa composante verticale ?
\item En déduire la direction, le sens et la valeur moyenne de $\dfrac{\Delta\vec{v}}{\Delta t}$.
\item La balle a une masse de \qty{60}{g}. Calculer la somme des forces et l'identifier.
\end{questions}
\end{exercice}
\end{blocLivre}
\begin{blocLivre}{ex-13-7}
\begin{exercice}[2]{La balle au bout d'un fil}
Une balle de \qty{200}{g} glisse sans frottement sur une table horizontale, retenue par un fil fixé en O : elle tourne à vitesse constante
sur un cercle de rayon $R=\qty{0,50}{m}$. On a relevé ses positions toutes les \qty{0,10}{s}.
\schema{13-ex07}
\begin{questions}
\item Mesurer $M_2M_3$ sur le document et en déduire la valeur de la vitesse.
\item Représenter $\vec{v}_2$ et $\vec{v}_3$ avec l'échelle \qty{1}{cm} pour \qty{1}{m.s^{-1}}.
\item Construire $\Delta\vec{v}_2$ et mesurer sa valeur. Quelle est sa direction ?
\item En déduire la direction, le sens et la valeur de la force exercée par le fil sur la balle.
\end{questions}
\end{exercice}
\end{blocLivre}
\begin{blocLivre}{ex-13-8}
\begin{exercice}[2]{Le tir au hockey}
Lors d'un tir, la crosse reste en contact avec le palet de \qty{170}{g}, initialement immobile, pendant \qty{10}{ms} ; le palet part à \qty{30}{m.s^{-1}}.
\begin{donnees}
$g=\qty{9,81}{N.kg^{-1}}$.
\end{donnees}
\begin{questions}
\item Calculer la variation de vitesse du palet pendant le contact.
\item Estimer la force moyenne exercée par la crosse sur le palet.
\item Comparer cette force au poids du palet. Peut-on négliger le poids pendant la frappe ?
\item Avec la même force, que deviendrait la vitesse du palet si le contact durait deux fois plus longtemps ?
\end{questions}
\end{exercice}
\end{blocLivre}
\begin{blocLivre}{ex-13-9}
\begin{exercice}[2]{Le saut en parachute}
\schema{13-ex09}
\begin{donnees}
$g=\qty{9,81}{N.kg^{-1}}$.
\end{donnees}
\begin{questions}
\item Au tout début du saut, seul le poids agit. Vérifier que la pente de la courbe est alors égale à $g$.
\item Entre \qty{25}{s} et \qty{40}{s}, que dire de la somme des forces ? En déduire la valeur des frottements de l'air.
\item À l'ouverture, la vitesse passe de \qty{55}{m.s^{-1}} à \qty{6}{m.s^{-1}} en \qty{3,0}{s}. Estimer la somme des forces, puis la force exercée par l'air.
\item Pourquoi le parachutiste atteint-il ensuite une nouvelle vitesse constante ?
\end{questions}
\end{exercice}
\end{blocLivre}
\begin{blocLivre}{ex-13-10}
\begin{exercice}[2]{Le chariot tiré}
Un chariot de \qty{500}{g}, tiré par un fil sur un rail horizontal, part du repos. On relève sa position toutes les \qty{0,10}{s}.
\par\smallskip\noindent\hfill\begin{tabular}{lccccccc}
\toprule
$t$ (s) & 0 & 0,10 & 0,20 & 0,30 & 0,40 & 0,50 & 0,60 \\
$x$ (m) & 0 & 0,006 & 0,024 & 0,054 & 0,096 & 0,150 & 0,216 \\
\bottomrule
\end{tabular}\hfill\null\par\smallskip
\begin{questions}
\item Calculer la vitesse moyenne du chariot sur chaque intervalle de \qty{0,10}{s}.
\item Calculer les variations de vitesse successives. Que remarque-t-on ?
\item En déduire la valeur de la somme des forces qui s'exercent sur le chariot.
\item Avec le même fil tirant de la même façon, comment varierait la vitesse d'un chariot deux fois plus lourd ?
\end{questions}
\end{exercice}
\end{blocLivre}
\begin{blocLivre}{ex-13-11}
\begin{exercice}[2]{Construire une variation de vitesse}
\schema{13-ex11}
\begin{questions}
\item Calculer les coordonnées des vecteurs vitesse $\vec{v}_2$ et $\vec{v}_3$, puis leurs valeurs.
\item Représenter $\vec{v}_2$ en $M_2$ et $\vec{v}_3$ en $M_3$ avec l'échelle \qty{1}{cm} pour \qty{2}{m.s^{-1}}.
\item Construire $\Delta\vec{v}_2$ en $M_3$ et déterminer sa valeur.
\item En déduire la somme des forces qui s'exercent sur la balle et la comparer à son poids ($g=\qty{9,81}{N.kg^{-1}}$).
\end{questions}
\end{exercice}
\end{blocLivre}
\rubrique{Objectif DS}
\begin{blocLivre}{ex-13-12}
\begin{exercice}[3]{Le service au tennis}
\begin{docu}{1}{La frappe}
Au service, le joueur lance la balle de \qty{58}{g} au-dessus de sa tête et la frappe quand elle est au sommet de sa trajectoire,
pratiquement immobile. La raquette reste en contact avec la balle pendant environ \qty{5,0}{ms}.
\end{docu}
\begin{docu}{2}{Des records}
Un bon joueur de club sert à \qty{180}{km/h}, soit \qty{50}{m.s^{-1}} ; les meilleurs professionnels dépassent \qty{220}{km/h}.
\end{docu}
\begin{donnees}
$g=\qty{9,81}{N.kg^{-1}}$.
\end{donnees}
\begin{questions}
\item Estimer la force moyenne exercée par la raquette sur la balle lors d'un service à \qty{180}{km/h}.
\item Comparer cette force au poids de la balle. Conclure.
\item Après la frappe, on néglige l'action de l'air. Quelle est la variation du vecteur vitesse de la balle en \qty{0,10}{s} ?
\item \initiative Estimer la force exercée par la raquette d'un professionnel servant à \qty{220}{km/h}, avec la même durée de contact.
\end{questions}
\end{exercice}
\end{blocLivre}
\begin{blocLivre}{ex-13-13}
\begin{exercice}[3]{Le décollage d'un lanceur}
\begin{docu}{1}{Au décollage}
Un lanceur lourd de \qty{780}{t} décolle verticalement. Ses moteurs exercent une poussée verticale de \qty{1,40e7}{N}, qu'on suppose
constante pendant les premières minutes du vol. On néglige l'action de l'air au départ.
\end{docu}
\begin{docu}{2}{Deux minutes plus tard}
Les moteurs ont brûlé une grande partie des ergols : la masse du lanceur n'est plus que de \qty{380}{t}.
\end{docu}
\begin{donnees}
$g=\qty{9,81}{N.kg^{-1}}$.
\end{donnees}
\begin{questions}
\item Faire l'inventaire des forces au décollage et les représenter à l'échelle \qty{1}{cm} pour \qty{4,0e6}{N}.
\item Calculer la somme des forces et en déduire la variation de vitesse pendant la première seconde.
\item \initiative Estimer la variation de vitesse en une seconde deux minutes plus tard, en négligeant toujours l'air.
\item Expliquer pourquoi la vitesse d'un lanceur augmente de plus en plus vite.
\end{questions}
\end{exercice}
\end{blocLivre}
\begin{blocLivre}{ex-13-14}
\begin{exercice}[3]{Ceinture et airbag}
\begin{docu}{1}{Un choc frontal}
Lors d'un choc contre un obstacle fixe à \qty{60}{km/h}, la voiture s'écrase et s'arrête en un dixième de seconde environ. Retenu par
sa ceinture et l'airbag, un passager de \qty{75}{kg} s'arrête en \qty{0,12}{s}. Sans protection, il heurterait le tableau de bord et s'arrêterait en
\qty{0,010}{s} environ.
\end{docu}
\begin{donnees}
$g=\qty{9,81}{N.kg^{-1}}$.
\end{donnees}
\begin{questions}
\item Calculer la variation de vitesse du passager, en \unit{m.s^{-1}}.
\item Estimer la force exercée sur le passager par la ceinture et l'airbag, puis la comparer à son poids.
\item Estimer la force exercée par le tableau de bord sur un passager non attaché.
\item \initiative Expliquer, à l'aide de la relation approchée, le rôle de la ceinture et de l'airbag.
\end{questions}
\end{exercice}
\end{blocLivre}
\begin{blocLivre}{ex-13-15}
\begin{exercice}[3]{Simuler un tir}
\begin{docu}{1}{Le principe}
Pour simuler un mouvement, on part de la position et de la vitesse initiales, puis on avance pas à pas : à chaque pas $\Delta t$, la
position change de $\vec{v}\,\Delta t$ et la vitesse de $\Delta\vec{v}=\dfrac{\sum\vec{F}}{m}\,\Delta t$.
\end{docu}
\begin{docu}{2}{Le programme}
\begin{python}
from math import cos, sin, radians
g, dt = 9.81, 0.01
vx, vy = 20 * cos(radians(45)), 20 * sin(radians(45))   # m/s
x, y = 0, 0
while y >= 0:
    x, y = x + vx * dt, y + vy * dt
    vy = vy - g * dt
print(round(x, 1), 'm')
\end{python}
\begin{console}
41.0 m
\end{console}
\end{docu}
\begin{questions}
\item \capacite{Capacité numérique}Quelles sont la vitesse et la direction initiales du projectile ? Que représente la valeur affichée ?
\item Quelle hypothèse sur les forces traduit la ligne \code{vy = vy - g * dt} ? Pourquoi \code{vx} n'est-il jamais modifié ?
\item Pourquoi la masse du projectile n'apparaît-elle pas dans le programme ?
\item \initiative Sur la Lune ($g=\qty{1,62}{N.kg^{-1}}$), le programme affiche \code{247.1 m}. Comparer et interpréter.
\end{questions}
\end{exercice}
\end{blocLivre}
